% The build command can override this default before reading the document.
\providecommand{\BayMode}{animated}
\documentclass[\BayMode]{baybeamer}

% Mark animation steps with \pause, \uncover<2->{...}, or [<+->].
\title{Programming with clarity}
\author[Your name]{Your full name} % Optional short name is used in the footer.
\date{\today}

\begin{document}
\begin{frame}[plain]
  \titlepage
\end{frame}

\section{Concepts}
\subsection{Definitions and examples}
\begin{frame}
  \begin{definitionbox}{Pure function}
    A pure function returns the same result for the same input and has
    no observable side effects.
  \end{definitionbox}
  \begin{examplebox}{Squaring a number}
    The function $f(x)=x^2$ maps $4$ to $16$, every time it is evaluated.
  \end{examplebox}
\end{frame}

\subsection{A step-by-step explanation}
\begin{frame}
  \begin{enumerate}[<+->]
    \item Start with the input values: $[1,2,3,4]$.
    \item Square each value: $[1,4,9,16]$.
    \item Add the results: $1+4+9+16=30$.
  \end{enumerate}
\end{frame}

\section{Code listings}
\subsection{Java: sum of squares}
% Frames containing literal code must be marked fragile.
\begin{frame}[fragile]
\begin{java}
public class SumOfSquares {
    static int sumOfSquares(int[] values) {
        int total = 0;
        for (int value : values) {
            total += value * value;
        }
        return total;
    }
}
\end{java}
\end{frame}

\subsection{Haskell: sum of squares}
\begin{frame}[fragile]
\begin{haskell}
-- Compose two pure functions.
sumOfSquares :: [Integer] -> Integer
sumOfSquares = sum . map (\x -> x * x)

main :: IO ()
main = print (sumOfSquares [1, 2, 3, 4])
\end{haskell}
  \begin{hintbox}{Function composition}
    Read the composition from right to left: first map, then sum.
  \end{hintbox}
\end{frame}

\subsection{Python: sum of squares}
\begin{frame}[fragile]
\begin{python}
def sum_of_squares(values: list[int]) -> int:
    # A generator avoids building an intermediate list.
    return sum(value * value for value in values)

print(sum_of_squares([1, 2, 3, 4]))  # 30
\end{python}
Inline identifiers use Fira Mono: \texttt{sum\_of\_squares}.
\end{frame}

\section{Practice}
\subsection{Exercise and solution}
\begin{frame}
  \begin{examplebox}{Try it yourself}
    Compute the sum of squares for $[2,3,5]$.
  \end{examplebox}
  \begin{hintbox}{Break it down}
    Square each number before adding the results.
  \end{hintbox}
  \pause
  \begin{solutionbox}{Worked answer}
    $2^2+3^2+5^2=4+9+25=38$.
  \end{solutionbox}
\end{frame}

\subsection{Takeaways}
\begin{frame}
  \begin{itemize}
    \item The same algorithm can be expressed in different languages.
    \item Definitions, examples, hints, and solutions have distinct boxes.
    \item The handout includes all steps and solutions on one page per frame.
  \end{itemize}
\end{frame}
\end{document}
